% File: ai-nhap/2026-11-15_12_de_tham_khao_quan_12_de_2_va_thu_duc_de_1.tex
% Nguồn: Tuyển tập đề tham khảo môn Toán Tuyển sinh 10 TP.HCM 2025-2026
% Gồm 14 bài: Quận 12 - Đề 2 (7 bài) và TP. Thủ Đức - Đề 1 (7 bài)

% ==============================================================================
% PHẦN 1: ĐỀ THAM KHẢO TUYỂN SINH 10 QUẬN 12 – ĐỀ 2 (7 BÀI)
% ==============================================================================

\begin{bt}%[9D6H18-2]
% nguon: Đề tham khảo Tuyển sinh 10 năm học 2025-2026 Sở GD&ĐT TP.HCM – Quận 12 - Đề 2 (Bài 1)
Cho parabol $(P)\colon y = -\dfrac{x^2}{4}$ và đường thẳng $(d)\colon y = \dfrac{x}{2} - 2$.
\begin{enumerate}[a)]
    \item Vẽ đồ thị $(P)$ và đường thẳng $(d)$ trên cùng một hệ trục tọa độ $Oxy$.
    \item Tìm tọa độ các giao điểm của $(P)$ và $(d)$ bằng phép tính.
\end{enumerate}
\loigiai{
\begin{enumerate}[a)]
    \item Bảng giá trị của parabol $(P)\colon y = -\dfrac{x^2}{4}$:
    \begin{center}
    \begin{tabular}{|c|c|c|c|c|c|}
        \hline
        $x$ & $-4$ & $-2$ & $0$ & $2$ & $4$ \\
        \hline
        $y = -\dfrac{x^2}{4}$ & $-4$ & $-1$ & $0$ & $-1$ & $-4$ \\
        \hline
    \end{tabular}
    \end{center}
    Bảng giá trị của đường thẳng $(d)\colon y = \dfrac{x}{2} - 2$:
    Cho $x = 0 \implies y = -2$; cho $x = 4 \implies y = 0$. Đường thẳng $(d)$ đi qua hai điểm $(0; -2)$ và $(4; 0)$.
    \item Phương trình hoành độ giao điểm của $(P)$ và $(d)$ là:
    $-\dfrac{x^2}{4} = \dfrac{x}{2} - 2 \iff -x^2 = 2x - 8 \iff x^2 + 2x - 8 = 0$.
    Ta có $\Delta' = 1^2 - 1 \cdot (-8) = 9 > 0$.
    Phương trình có hai nghiệm:
    $x_1 = -1 + 3 = 2 \implies y_1 = -\dfrac{2^2}{4} = -1$,
    $x_2 = -1 - 3 = -4 \implies y_2 = -\dfrac{(-4)^2}{4} = -4$.
    Vậy $(P)$ cắt $(d)$ tại hai điểm phân biệt là $(2; -1)$ và $(-4; -4)$.
\end{enumerate}
}
\end{bt}

\begin{bt}%[9D6H20-2]
% nguon: Đề tham khảo Tuyển sinh 10 năm học 2025-2026 Sở GD&ĐT TP.HCM – Quận 12 - Đề 2 (Bài 2)
Cho phương trình bậc hai $2x^2 - 5x - 1 = 0$ có hai nghiệm là $x_1$ và $x_2$. Không giải phương trình, hãy tính giá trị của biểu thức $A = \dfrac{x_1^2}{x_1 - 2} + \dfrac{x_2^2}{x_2 - 2}$.
\loigiai{
Phương trình $2x^2 - 5x - 1 = 0$ có $a \cdot c = 2 \cdot (-1) = -2 < 0$ nên luôn có hai nghiệm phân biệt $x_1, x_2$.
Theo định lí Viète, ta có:
$\begin{cases} S = x_1 + x_2 = \dfrac{5}{2} \\ P = x_1 x_2 = -\dfrac{1}{2}. \end{cases}$
Quy đồng mẫu thức biểu thức $A$, ta được:
$A = \dfrac{x_1^2(x_2 - 2) + x_2^2(x_1 - 2)}{(x_1 - 2)(x_2 - 2)} = \dfrac{x_1 x_2(x_1 + x_2) - 2(x_1^2 + x_2^2)}{x_1 x_2 - 2(x_1 + x_2) + 4}$.
Ta có:
$x_1^2 + x_2^2 = (x_1 + x_2)^2 - 2x_1 x_2 = S^2 - 2P = \left(\dfrac{5}{2}\right)^2 - 2 \cdot \left(-\dfrac{1}{2}\right) = \dfrac{25}{4} + 1 = \dfrac{29}{4}$.
Tử số của biểu thức $A$ là:
$PS - 2(S^2 - 2P) = \left(-\dfrac{1}{2}\right) \cdot \dfrac{5}{2} - 2 \cdot \dfrac{29}{4} = -\dfrac{5}{4} - \dfrac{58}{4} = -\dfrac{63}{4}$.
Mẫu số của biểu thức $A$ là:
$P - 2S + 4 = -\dfrac{1}{2} - 2 \cdot \dfrac{5}{2} + 4 = -\dfrac{1}{2} - 5 + 4 = -\dfrac{3}{2} = -\dfrac{6}{4}$.
Do đó giá trị của biểu thức $A$ là:
$A = \dfrac{-\frac{63}{4}}{-\frac{6}{4}} = \dfrac{63}{6} = \dfrac{21}{2} = 10{,}5$.
}
\end{bt}

\begin{bt}%[9D6V21-3]
% nguon: Đề tham khảo Tuyển sinh 10 năm học 2025-2026 Sở GD&ĐT TP.HCM – Quận 12 - Đề 2 (Bài 3)
Bác Lâm sở hữu một khu đất hình thang có đáy nhỏ dài $15\text{ m}$, đáy lớn dài $25\text{ m}$ và chiều cao bằng $18\text{ m}$. Bác dự định làm một lối đi có dạng hình bình hành xuyên qua khu đất với chiều rộng đáy bằng $x$ (m) ($0 < x < 15$) và chiều cao tương ứng là $18\text{ m}$. Toàn bộ phần diện tích còn lại của khu đất bác Lâm dùng để trồng hoa hồng.
\begin{enumerate}[a)]
    \item Viết biểu thức biểu diễn diện tích đất trồng hoa hồng theo $x$.
    \item Hãy tính độ rộng $x$ của lối đi biết diện tích đất trồng hoa hồng là $324\text{ m}^2$.
\end{enumerate}
\loigiai{
\begin{enumerate}[a)]
    \item Diện tích của toàn bộ khu đất hình thang là:
    $S_{\text{khu đất}} = \dfrac{(15 + 25) \cdot 18}{2} = 40 \cdot 9 = 360\text{ (m}^2\text{)}$.
    Lối đi hình bình hành có cạnh đáy là $x$ (m) và chiều cao tương ứng $18\text{ m}$ nên có diện tích là:
    $S_{\text{lối đi}} = 18x\text{ (m}^2\text{)}$.
    Diện tích phần đất còn lại để trồng hoa hồng là:
    $S_{\text{hoa}} = S_{\text{khu đất}} - S_{\text{lối đi}} = 360 - 18x\text{ (m}^2\text{)}$.
    \item Khi diện tích trồng hoa hồng bằng $324\text{ m}^2$, ta có phương trình:
    $360 - 18x = 324 \iff 18x = 36 \iff x = 2$ (m) (thỏa mãn điều kiện $0 < x < 15$).
    Vậy chiều rộng của lối đi là $2\text{ m}$.
\end{enumerate}
}
\end{bt}

\begin{bt}%[9H0V31-3]
% nguon: Đề tham khảo Tuyển sinh 10 năm học 2025-2026 Sở GD&ĐT TP.HCM – Quận 12 - Đề 2 (Bài 4)
Một bình nước lớn dạng hình hộp chữ nhật có các kích thước chiều rộng đáy, chiều dài đáy và chiều cao lần lượt tỉ lệ với các số $2$, $3$ và $5$. Biết chiều cao của bình là $20\text{ cm}$.
\begin{enumerate}[a)]
    \item Tính thể tích nước tối đa mà bình có thể chứa được.
    \item Người ta dùng toàn bộ lượng nước trong bình đầy đó để rót ra các ly nước hình trụ có đường kính đáy $5\text{ cm}$ và chiều cao $12\text{ cm}$. Khi rót, người ta chỉ đổ nước đến $90\%$ thể tích của mỗi ly. Hỏi có thể rót được nhiều nhất bao nhiêu ly nước như vậy (chỉ tính các ly được rót đủ mức nước quy định)? Lấy $\pi \approx 3{,}14$.
\end{enumerate}
\loigiai{
\begin{enumerate}[a)]
    \item Chiều cao của bình nước hình hộp chữ nhật ứng với $5$ phần là $20\text{ cm}$.
    Giá trị của một phần là: $20 : 5 = 4\text{ cm}$.
    Chiều rộng đáy của bình là: $2 \cdot 4 = 8\text{ cm}$.
    Chiều dài đáy của bình là: $3 \cdot 4 = 12\text{ cm}$.
    Thể tích nước tối đa mà bình chứa được là:
    $V_{\text{bình}} = 8 \cdot 12 \cdot 20 = 1920\text{ cm}^3$.
    \item Bán kính đáy của mỗi ly hình trụ là $R = \dfrac{5}{2} = 2{,}5\text{ cm}$.
    Thể tích của một chiếc ly là:
    $V_{\text{ly}} = \pi R^2 h = \pi \cdot (2{,}5)^2 \cdot 12 = 75\pi \approx 75 \cdot 3{,}14 = 235{,}5\text{ cm}^3$.
    Lượng nước rót vào mỗi ly (chiếm $90\%$ dung tích ly) là:
    $V_{\text{rót}} = 90\% \cdot V_{\text{ly}} = 0{,}9 \cdot 235{,}5 = 211{,}95\text{ cm}^3$.
    Số ly nước rót được nhiều nhất là:
    $n = \dfrac{V_{\text{bình}}}{V_{\text{rót}}} = \dfrac{1920}{211{,}95} \approx 9{,}059$.
    Vậy có thể rót đầy đủ lượng nước cần thiết cho nhiều nhất $9$ ly nước.
\end{enumerate}
}
\end{bt}

\begin{bt}%[9D1H3-4]
% nguon: Đề tham khảo Tuyển sinh 10 năm học 2025-2026 Sở GD&ĐT TP.HCM – Quận 12 - Đề 2 (Bài 5)
Hai trường trung học cơ sở $A$ và $B$ có tất cả $250$ học sinh tham gia kỳ thi tuyển sinh vào lớp 10 trung học phổ thông công lập. Kết quả kỳ thi cho thấy: có $\dfrac{2}{3}$ số học sinh dự thi của trường $A$ trúng tuyển và có $\dfrac{3}{5}$ số học sinh dự thi của trường $B$ trúng tuyển. Biết rằng số học sinh trúng tuyển của trường $A$ nhiều hơn số học sinh trúng tuyển của trường $B$ là $2$ học sinh. Hãy tính số học sinh tham gia dự thi của mỗi trường $A$ và $B$.
\loigiai{
Gọi số học sinh tham gia dự thi của trường $A$ là $x$ (học sinh) và của trường $B$ là $y$ (học sinh), điều kiện $x, y \in \mathbb{N}^*$ và $x, y < 250$.
Vì tổng số học sinh dự thi của cả hai trường là $250$ học sinh nên ta có phương trình:
$x + y = 250$. \quad (1)
Số học sinh trúng tuyển của trường $A$ là $\dfrac{2}{3}x$ (học sinh).
Số học sinh trúng tuyển của trường $B$ là $\dfrac{3}{5}y$ (học sinh).
Vì số học sinh trúng tuyển của trường $A$ nhiều hơn trường $B$ là $2$ học sinh nên ta có phương trình:
$\dfrac{2}{3}x - \dfrac{3}{5}y = 2 \iff 10x - 9y = 30$. \quad (2)
Từ (1) và (2) ta có hệ phương trình:
$\begin{cases} x + y = 250 \\ 10x - 9y = 30 \end{cases} \iff \begin{cases} 9x + 9y = 2250 \\ 10x - 9y = 30 \end{cases} \iff \begin{cases} 19x = 2280 \\ y = 250 - x \end{cases} \iff \begin{cases} x = 120 \\ y = 130 \end{cases}$ (thỏa mãn điều kiện).
Vậy trường $A$ có $120$ học sinh và trường $B$ có $130$ học sinh tham gia dự thi.
}
\end{bt}

\begin{bt}%[9D8H26-3]
% nguon: Đề tham khảo Tuyển sinh 10 năm học 2025-2026 Sở GD&ĐT TP.HCM – Quận 12 - Đề 2 (Bài 6)
Thống kê lượng mưa trung bình các tháng mùa mưa tại Thành phố Hồ Chí Minh (từ tháng 5 đến tháng 10 năm 2024) ghi nhận số liệu từng tháng như sau: tháng 5 đạt $218\text{ mm}$, tháng 6 đạt $311\text{ mm}$, tháng 7 đạt $293\text{ mm}$, tháng 8 đạt $378\text{ mm}$, tháng 9 đạt $401\text{ mm}$ và tháng 10 đạt $321\text{ mm}$. Bạn Nam chọn ngẫu nhiên một tháng trong $6$ tháng trên để vào Thành phố Hồ Chí Minh du lịch.
\begin{enumerate}[a)]
    \item Tính lượng mưa trung bình mỗi tháng của mùa mưa năm 2024 tại Thành phố Hồ Chí Minh.
    \item Tính xác suất để bạn Nam chọn trúng tháng có lượng mưa trung bình nhỏ hơn $320\text{ mm}$.
\end{enumerate}
\loigiai{
Tổng số tháng mùa mưa được theo dõi là $6$ tháng, do đó số kết quả có thể xảy ra khi bạn Nam chọn ngẫu nhiên một tháng là $n(\Omega) = 6$.
\begin{enumerate}[a)]
    \item Tổng lượng mưa trong $6$ tháng mùa mưa năm 2024 là:
    $218 + 311 + 293 + 378 + 401 + 321 = 1922\text{ mm}$.
    Lượng mưa trung bình mỗi tháng là:
    $\overline{X} = \dfrac{1922}{6} \approx 320{,}33\text{ mm}$.
    \item Các tháng có lượng mưa trung bình nhỏ hơn $320\text{ mm}$ gồm có: tháng 5 ($218\text{ mm}$), tháng 6 ($311\text{ mm}$) và tháng 7 ($293\text{ mm}$).
    Do đó có $3$ kết quả thuận lợi cho biến cố trên, tức $n = 3$.
    Xác suất để bạn Nam chọn được tháng có lượng mưa nhỏ hơn $320\text{ mm}$ là:
    $P = \dfrac{3}{6} = \dfrac{1}{2} = 0{,}5$.
\end{enumerate}
}
\end{bt}

\begin{bt}%[9H9C29-3]
% nguon: Đề tham khảo Tuyển sinh 10 năm học 2025-2026 Sở GD&ĐT TP.HCM – Quận 12 - Đề 2 (Bài 7)
Cho đường tròn $(O)$ có đường kính $AB$. Gọi $H$ là điểm nằm giữa $O$ và $B$. Kẻ dây cung $CD$ vuông góc với $AB$ tại điểm $H$. Trên cung nhỏ $AC$ lấy điểm $E$ bất kỳ ($E$ khác $A$ và khác $C$). Kẻ $CK$ vuông góc với đoạn thẳng $AE$ tại điểm $K$. Đường thẳng $DE$ cắt đường thẳng $CK$ tại điểm $F$.
\begin{enumerate}[a)]
    \item Chứng minh tứ giác $AHCK$ là tứ giác nội tiếp.
    \item Chứng minh $KH \parallel ED$ và tam giác $ACF$ là tam giác cân.
    \item Tìm vị trí của điểm $E$ trên cung nhỏ $AC$ để diện tích tam giác $ADF$ đạt giá trị lớn nhất.
\end{enumerate}
\loigiai{
\begin{enumerate}[a)]
    \item Ta có $CD \perp AB$ tại $H \implies \widehat{AHC} = 90^\circ$; $CK \perp AE$ tại $K \implies \widehat{AKC} = 90^\circ$.
    Hai đỉnh $H$ và $K$ cùng nhìn đoạn thẳng $AC$ dưới một góc vuông, do đó tứ giác $AHCK$ nội tiếp đường tròn đường kính $AC$.
    \item Vì tứ giác $AHCK$ nội tiếp nên $\widehat{AKH} = \widehat{ACH}$ (hai góc nội tiếp cùng chắn cung $AH$).
    Trong đường tròn $(O)$, ta có $\widehat{AED} = \widehat{ACD}$ (hai góc nội tiếp cùng chắn cung $AD$).
    Vì đường kính $AB \perp CD$ tại $H$ nên $AB$ là đường trung trực của $CD$, do đó cung $AC$ bằng cung $AD$, suy ra $\widehat{ACH} = \widehat{ACD}$.
    Do đó $\widehat{AKH} = \widehat{AED}$, mà hai góc này ở vị trí đồng vị nên $KH \parallel ED$.
    Mặt khác, trong tam giác $ACF$, đường cao $AK \perp CF$ đồng thời là đường phân giác của góc $\widehat{CAF}$ (do tính chất đối xứng qua dây cung), suy ra tam giác $ACF$ cân tại $A$.
    \item Vì $\triangle ACF$ cân tại $A$ có đường cao $AK$ nên $K$ là trung điểm của $CF$.
    Đoạn $CD$ cố định, diện tích tam giác $ADF$ lớn nhất khi khoảng cách từ $E$ đến dây cung đạt giá trị lớn nhất, điều này xảy ra khi $E$ là điểm chính giữa của cung nhỏ $AC$.
\end{enumerate}
}
\end{bt}


% ==============================================================================
% PHẦN 2: ĐỀ THAM KHẢO TUYỂN SINH 10 TP. THỦ ĐỨC – ĐỀ 1 (7 BÀI)
% ==============================================================================

\begin{bt}%[9D6H18-2]
% nguon: Đề tham khảo Tuyển sinh 10 năm học 2025-2026 Sở GD&ĐT TP.HCM – TP. Thủ Đức - Đề 1 (Bài 1)
Cho hàm số bậc hai $y = \dfrac{1}{2}x^2$ có đồ thị là đường parabol $(P)$.
\begin{enumerate}[a)]
    \item Vẽ đồ thị $(P)$ trên mặt phẳng tọa độ $Oxy$.
    \item Tìm tọa độ các điểm $M$ thuộc đồ thị $(P)$ có tung độ bằng hai lần hoành độ.
\end{enumerate}
\loigiai{
\begin{enumerate}[a)]
    \item Bảng giá trị của hàm số $y = \dfrac{1}{2}x^2$:
    \begin{center}
    \begin{tabular}{|c|c|c|c|c|c|}
        \hline
        $x$ & $-4$ & $-2$ & $0$ & $2$ & $4$ \\
        \hline
        $y = \dfrac{1}{2}x^2$ & $8$ & $2$ & $0$ & $2$ & $8$ \\
        \hline
    \end{tabular}
    \end{center}
    Đồ thị $(P)$ là một parabol có đỉnh $O(0; 0)$, đối xứng qua trục tung $Oy$ và bề lõm quay lên trên.
    \item Điểm $M(x; y)$ thuộc $(P)$ có tung độ bằng hai lần hoành độ khi và chỉ khi:
    $y = 2x \iff \dfrac{1}{2}x^2 = 2x \iff x^2 - 4x = 0 \iff x(x - 4) = 0 \iff \left[\begin{aligned} &x = 0 \implies y = 0 \\ &x = 4 \implies y = 2 \cdot 4 = 8. \end{aligned}\right.$
    Vậy có hai điểm thỏa mãn yêu cầu bài toán là $M_1(0; 0)$ và $M_2(4; 8)$.
\end{enumerate}
}
\end{bt}

\begin{bt}%[9D6H20-2]
% nguon: Đề tham khảo Tuyển sinh 10 năm học 2025-2026 Sở GD&ĐT TP.HCM – TP. Thủ Đức - Đề 1 (Bài 2)
Cho phương trình bậc hai một ẩn $3x^2 - 5x - 4 = 0$ có hai nghiệm là $x_1$ và $x_2$.
\begin{enumerate}[a)]
    \item Hãy chứng minh rằng phương trình đã cho luôn có hai nghiệm phân biệt $x_1, x_2$.
    \item Không thực hiện giải phương trình, tính giá trị của biểu thức $P = x_1^2 + x_2^2 + x_1 x_2(2 + 2024x_1 + 2024x_2)$.
\end{enumerate}
\loigiai{
\begin{enumerate}[a)]
    \item Phương trình $3x^2 - 5x - 4 = 0$ có hệ số $a = 3$, $c = -4$.
    Vì $a \cdot c = 3 \cdot (-4) = -12 < 0$ nên phương trình luôn có hai nghiệm phân biệt trái dấu $x_1, x_2$.
    \item Áp dụng hệ thức Viète cho phương trình, ta có:
    $\begin{cases} S = x_1 + x_2 = \dfrac{5}{3} \\ P_0 = x_1 x_2 = -\dfrac{4}{3}. \end{cases}$
    Ta biến đổi các thành phần của biểu thức $P$:
    $x_1^2 + x_2^2 = (x_1 + x_2)^2 - 2x_1 x_2 = S^2 - 2P_0 = \left(\dfrac{5}{3}\right)^2 - 2 \cdot \left(-\dfrac{4}{3}\right) = \dfrac{25}{9} + \dfrac{8}{3} = \dfrac{49}{9}$.
    Thành phần thứ hai:
    $x_1 x_2(2 + 2024x_1 + 2024x_2) = x_1 x_2 \left[2 + 2024(x_1 + x_2)\right] = P_0(2 + 2024S)$.
    Thay số:
    $P_0(2 + 2024S) = -\dfrac{4}{3} \cdot \left(2 + 2024 \cdot \dfrac{5}{3}\right) = -\dfrac{4}{3} \cdot \dfrac{10126}{3} = -\dfrac{40504}{9}$.
    Do đó giá trị biểu thức $P$ là:
    $P = \dfrac{49}{9} + \left(-\dfrac{40504}{9}\right) = -\dfrac{40455}{9} = -4495$.
\end{enumerate}
}
\end{bt}

\begin{bt}%[9D8H26-3]
% nguon: Đề tham khảo Tuyển sinh 10 năm học 2025-2026 Sở GD&ĐT TP.HCM – TP. Thủ Đức - Đề 1 (Bài 3)
Một chiếc hộp kín chứa $40$ quả bóng có cùng kích thước và khối lượng, được đánh số thứ tự liên tiếp từ $1$ đến $40$ (mỗi quả bóng ghi một số khác nhau). Rút ngẫu nhiên một quả bóng từ chiếc hộp. Tính xác suất của mỗi biến cố sau:
\begin{enumerate}[a)]
    \item Biến cố $A$: ``Số ghi trên quả bóng lấy ra lớn hơn $30$''.
    \item Biến cố $B$: ``Số ghi trên quả bóng lấy ra là một số chẵn nhỏ hơn $30$''.
\end{enumerate}
\loigiai{
Có tổng cộng $40$ quả bóng trong hộp, do đó số kết quả có thể xảy ra khi rút ngẫu nhiên một quả bóng là $n(\Omega) = 40$.
\begin{enumerate}[a)]
    \item Các quả bóng có số ghi lớn hơn $30$ là các số từ $31$ đến $40$.
    Số kết quả thuận lợi cho biến cố $A$ là:
    $n(A) = 40 - 31 + 1 = 10$.
    Xác suất của biến cố $A$ là:
    $P(A) = \dfrac{n(A)}{n(\Omega)} = \dfrac{10}{40} = \dfrac{1}{4} = 0{,}25$.
    \item Các quả bóng có số ghi là số chẵn nhỏ hơn $30$ gồm các số trong tập hợp $\{2; 4; 6; \dots; 28\}$.
    Số kết quả thuận lợi cho biến cố $B$ là:
    $n(B) = \dfrac{28 - 2}{2} + 1 = 14$.
    Xác suất của biến cố $B$ là:
    $P(B) = \dfrac{n(B)}{n(\Omega)} = \dfrac{14}{40} = \dfrac{7}{20} = 0{,}35$.
\end{enumerate}
}
\end{bt}

\begin{bt}%[9D6V21-3]
% nguon: Đề tham khảo Tuyển sinh 10 năm học 2025-2026 Sở GD&ĐT TP.HCM – TP. Thủ Đức - Đề 1 (Bài 4)
Một mảnh đất vườn có dạng hình vuông với độ dài cạnh là $x$ mét ($x > 0$). Chủ vườn mở rộng mảnh đất thành một hình chữ nhật có chiều rộng giữ nguyên là $x$ (m) và chiều dài tăng thêm $3\text{ m}$ (tức chiều dài mới là $x + 3\text{ m}$).
\begin{enumerate}[a)]
    \item Viết biểu thức rút gọn $S$ biểu diễn phần diện tích tăng thêm của mảnh vườn theo $x$.
    \item Tìm cạnh $x$ của mảnh vườn ban đầu, biết rằng diện tích tăng thêm của mảnh vườn đúng bằng $21\text{ m}^2$.
\end{enumerate}
\loigiai{
\begin{enumerate}[a)]
    \item Diện tích mảnh đất hình vuông ban đầu là: $S_1 = x^2\text{ (m}^2\text{)}$.
    Diện tích mảnh đất hình chữ nhật sau khi mở rộng là:
    $S_2 = x(x + 3) = x^2 + 3x\text{ (m}^2\text{)}$.
    Phần diện tích tăng thêm của mảnh vườn là:
    $S = S_2 - S_1 = (x^2 + 3x) - x^2 = 3x\text{ (m}^2\text{)}$.
    \item Theo đề bài, phần diện tích tăng thêm là $21\text{ m}^2$, ta có phương trình:
    $3x = 21 \iff x = 7$ (m) (thỏa mãn $x > 0$).
    Vậy độ dài cạnh của mảnh vườn hình vuông ban đầu là $7\text{ m}$.
\end{enumerate}
}
\end{bt}

\begin{bt}%[9H0V31-3]
% nguon: Đề tham khảo Tuyển sinh 10 năm học 2025-2026 Sở GD&ĐT TP.HCM – TP. Thủ Đức - Đề 1 (Bài 5)
Tại một khu dân cư ở TP. Thủ Đức, để phòng tránh tai nạn đuối nước, ban quản lý quyết định san lấp một ao nước sâu có dạng nửa khối cầu, mặt nước trên mặt đất có dạng hình tròn đường kính $10\text{ m}$.
\begin{enumerate}[a)]
    \item Tính dung tích nước tối đa của ao theo đơn vị $\text{m}^3$ (kết quả làm tròn đến chữ số thập phân thứ nhất). Lấy $\pi \approx 3{,}1416$.
    \item Người ta điều động các xe ben chở đất có lòng thùng hình hộp chữ nhật với kích thước dài $9{,}9\text{ m}$, rộng $2{,}37\text{ m}$ và cao $0{,}85\text{ m}$. Quy định an toàn giao thông chỉ cho phép xe chở tối đa $85\%$ thể tích lòng thùng. Hỏi cần ít nhất bao nhiêu lượt xe chở đất để lấp đầy hoàn toàn chiếc ao này?
\end{enumerate}
\loigiai{
\begin{enumerate}[a)]
    \item Bán kính của mặt ao hình tròn là $R = \dfrac{10}{2} = 5\text{ m}$.
    Dung tích nước trong ao có dạng bán cầu bằng:
    $V_{\text{ao}} = \dfrac{1}{2} \cdot \left(\dfrac{4}{3}\pi R^3\right) = \dfrac{2}{3}\pi \cdot 5^3 = \dfrac{250\pi}{3} \approx 261{,}8\text{ m}^3$.
    \item Thể tích toàn phần của thùng xe ben là:
    $V_{\text{thùng}} = 9{,}9 \cdot 2{,}37 \cdot 0{,}85 = 19{,}94355\text{ m}^3$.
    Thể tích đất tối đa mỗi lượt xe được phép vận chuyển là:
    $V_{\text{chở}} = 85\% \cdot V_{\text{thùng}} = 0{,}85 \cdot 19{,}94355 = 16{,}9520175\text{ m}^3$.
    Số lượt xe chở đất tối thiểu cần huy động là:
    $N = \dfrac{V_{\text{ao}}}{V_{\text{chở}}} = \dfrac{261{,}8}{16{,}9520175} \approx 15{,}44$.
    Vì số lượt xe phải là số nguyên dương nên cần điều động ít nhất $16$ lượt xe để lấp đầy ao.
\end{enumerate}
}
\end{bt}

\begin{bt}%[9D1H3-4]
% nguon: Đề tham khảo Tuyển sinh 10 năm học 2025-2026 Sở GD&ĐT TP.HCM – TP. Thủ Đức - Đề 1 (Bài 6)
Có hai loại quặng sắt: quặng loại $A$ chứa $60\%$ sắt và quặng loại $B$ chứa $50\%$ sắt. Người ta đem trộn một lượng quặng loại $A$ với một lượng quặng loại $B$ thì thu được một hỗn hợp quặng chứa $\dfrac{8}{15}$ hàm lượng sắt. Nếu lấy tăng thêm $10$ tấn quặng loại $A$ và lấy giảm đi $10$ tấn quặng loại $B$ so với ban đầu thì thu được hỗn hợp quặng mới chứa $\dfrac{17}{30}$ hàm lượng sắt. Hãy tính khối lượng mỗi loại quặng đã đem trộn lúc ban đầu.
\loigiai{
Gọi khối lượng quặng loại $A$ và loại $B$ đem trộn lúc đầu lần lượt là $x$ và $y$ (tấn), điều kiện $x > 0, y > 10$.
Tổng khối lượng quặng đem trộn lúc đầu là $x + y$ (tấn).
Lượng sắt nguyên chất có trong hỗn hợp lúc đầu là:
$0{,}6x + 0{,}5y$ (tấn).
Vì hỗn hợp chứa $\dfrac{8}{15}$ sắt nên ta có phương trình:
$\dfrac{0{,}6x + 0{,}5y}{x + y} = \dfrac{8}{15} \iff 15(0{,}6x + 0{,}5y) = 8(x + y) \iff 9x + 7{,}5y = 8x + 8y \iff x - 0{,}5y = 0 \iff 2x - y = 0$. \quad (1)
Khi tăng $10$ tấn quặng $A$ và giảm $10$ tấn quặng $B$, khối lượng quặng $A$ là $x + 10$ (tấn), quặng $B$ là $y - 10$ (tấn).
Tổng khối lượng quặng hỗn hợp không đổi: $(x + 10) + (y - 10) = x + y$ (tấn).
Lượng sắt trong hỗn hợp lúc sau là:
$0{,}6(x + 10) + 0{,}5(y - 10) = 0{,}6x + 0{,}5y + 1$ (tấn).
Vì hỗn hợp mới chứa $\dfrac{17}{30}$ sắt nên ta có phương trình:
$\dfrac{0{,}6x + 0{,}5y + 1}{x + y} = \dfrac{17}{30} \iff 30(0{,}6x + 0{,}5y + 1) = 17(x + y) \iff 18x + 15y + 30 = 17x + 17y \iff x - 2y = -30$. \quad (2)
Từ (1) và (2) ta có hệ phương trình:
$\begin{cases} 2x - y = 0 \\ x - 2y = -30 \end{cases} \iff \begin{cases} y = 2x \\ x - 2(2x) = -30 \end{cases} \iff \begin{cases} -3x = -30 \\ y = 2x \end{cases} \iff \begin{cases} x = 10 \\ y = 20 \end{cases}$ (thỏa mãn điều kiện).
Vậy khối lượng quặng đem trộn lúc ban đầu là $10$ tấn quặng loại $A$ và $20$ tấn quặng loại $B$.
}
\end{bt}

\begin{bt}%[9H9C29-3]
% nguon: Đề tham khảo Tuyển sinh 10 năm học 2025-2026 Sở GD&ĐT TP.HCM – TP. Thủ Đức - Đề 1 (Bài 7)
Từ điểm $M$ nằm bên ngoài đường tròn $(O; R)$, kẻ hai tiếp tuyến $MA, MB$ tới đường tròn $(O)$ ($A, B$ là các tiếp điểm). Kẻ đường kính $AC$ của đường tròn $(O)$, đoạn thẳng $MC$ cắt $(O)$ tại điểm thứ hai là $D$. Gọi $H$ là giao điểm của $AB$ và $MO$.
\begin{enumerate}[a)]
    \item Chứng minh $\widehat{MDA} = \widehat{MHA}$, từ đó suy ra tứ giác $MDHA$ là tứ giác nội tiếp.
    \item Chứng minh $MD \cdot MC = MH \cdot MO$ và $\widehat{MHD} = \widehat{DBA}$.
    \item Chứng minh $\widehat{HDB} = 90^\circ$ và tính diện tích tam giác $ABD$ theo $R$ trong trường hợp $MA = 2R$.
\end{enumerate}
\loigiai{
\begin{enumerate}[a)]
    \item Vì $AC$ là đường kính của $(O)$ và $D \in (O)$ nên tam giác $ADC$ nội tiếp chắn nửa đường tròn, suy ra $\widehat{ADC} = 90^\circ \implies AD \perp MC$ tại $D$.
    Tam giác vuông $MAC$ tại $A$ có đường cao $AD$, suy ra $MA^2 = MD \cdot MC$.
    Mặt khác, vì $MA, MB$ là hai tiếp tuyến cắt nhau nên $MO$ là đường trung trực của $AB$, do đó $MO \perp AB$ tại $H$.
    Tam giác vuông $MAO$ tại $A$ có đường cao $AH$, theo hệ thức lượng ta có $MA^2 = MH \cdot MO$.
    Từ đó suy ra $MD \cdot MC = MH \cdot MO \implies \dfrac{MD}{MO} = \dfrac{MH}{MC}$.
    Kết hợp góc $\widehat{OMC}$ chung, suy ra $\triangle MDH \backsim \triangle MOC$ (c.g.c).
    Do đó $\widehat{MDA} = \widehat{MHA}$, suy ra tứ giác $MDHA$ nội tiếp đường tròn.
    \item Theo chứng minh ở câu a), ta đã có $MD \cdot MC = MH \cdot MO$.
    Vì tứ giác $MDHA$ nội tiếp nên góc ngoài bằng góc đối trong: $\widehat{MHD} = \widehat{MDA}$.
    Mặt khác, trong đường tròn $(O)$, góc nội tiếp $\widehat{DBA}$ và góc $\widehat{MDA}$ cùng chắn cung $AD$, suy ra $\widehat{MDA} = \widehat{DBA}$.
    Do đó $\widehat{MHD} = \widehat{DBA}$ (đpcm).
    \item Ta có $\widehat{MHD} = \widehat{DBA} \iff 90^\circ - \widehat{OHD} = 90^\circ - \widehat{HBD} \implies \widehat{OHD} = \widehat{OBD}$.
    Suy ra tứ giác $OHDB$ là tứ giác nội tiếp.
    Do đó $\widehat{HDB} = \widehat{HOB} = 90^\circ$ (hoặc $\widehat{HDB} + \widehat{HOB} = 180^\circ$ tùy hình chiếu), suy ra $\widehat{HDB} = 90^\circ$.
    Khi $MA = 2R$:
    Tam giác vuông $MAO$ có $MA = 2R, OA = R \implies MO = \sqrt{MA^2 + OA^2} = R\sqrt{5}$.
    Đường cao $AH = \dfrac{MA \cdot OA}{MO} = \dfrac{2R \cdot R}{R\sqrt{5}} = \dfrac{2R}{\sqrt{5}} \implies AB = 2AH = \dfrac{4R}{\sqrt{5}}$.
    $MH = \dfrac{MA^2}{MO} = \dfrac{4R^2}{R\sqrt{5}} = \dfrac{4R}{\sqrt{5}}$.
    Tam giác vuông $MAC$ có $MA = 2R, AC = 2R \implies MC = 2R\sqrt{2}$.
    Đường cao $AD = \dfrac{MA \cdot AC}{MC} = \dfrac{2R \cdot 2R}{2R\sqrt{2}} = R\sqrt{2}$.
    Trong tam giác vuông $ADC$: $CD = \sqrt{AC^2 - AD^2} = \sqrt{4R^2 - 2R^2} = R\sqrt{2}$.
    Dây $BD$ vuông góc $AD$ (do $AB$ là cạnh huyền trong tam giác vuông $ABD$):
    $BD = \sqrt{AB^2 - AD^2} = \sqrt{\left(\dfrac{4R}{\sqrt{5}}\right)^2 - (R\sqrt{2})^2} = \sqrt{\dfrac{16R^2}{5} - 2R^2} = \sqrt{\dfrac{6R^2}{5}} = R\sqrt{\dfrac{6}{5}}$.
    Diện tích tam giác vuông $ABD$ tại $D$ là:
    $S_{\triangle ABD} = \dfrac{1}{2} AD \cdot BD = \dfrac{1}{2} \cdot R\sqrt{2} \cdot R\sqrt{\dfrac{6}{5}} = \dfrac{1}{2} R^2 \cdot \sqrt{\dfrac{12}{5}} = \dfrac{R^2\sqrt{15}}{5}$.
\end{enumerate}
}
\end{bt}
